
\subsection{Introducción a Octave}

% ___________________________________________________________________________

\hypertarget{qu-es-octave-i}{}
\pdfbookmark[0]{Qué es Octave}{qu-es-octave-i}
\begin{frame}\frametitle{Qué es Octave}
\begin{itemize}
\item {} 
  \textbf{Lenguaje} interpretado de alto nivel, orientado al \textbf{cálculo matricial}

\item {} 
  \textbf{Intérprete} para este lenguaje

\item {} 
Licencia libre (GPL)
\begin{itemize}
\item {} 
Concebido en 1988 por John W. Eaton

\item {} 
Amplia comunidad de usuarios (\href{http://www.octave.org}{http://www.octave.org})

\item {} 
Código fuente C++
\begin{itemize}
\item {} 
Potente biblioteca de cálculo matricial (GPL)

\item {} 
Fácilmente extensible en C++

\end{itemize}

\item {} 
Disponible para numerosos sistemas y arquitecturas

\end{itemize}


\item Octave \emph{\bf no} está especializado en \emph{\bf cálculo simbólico}
    \begin{itemize}
    \item {} Para ello, existen numerosas herramientas libres como
       \textbf{Python}+\textbf{SymPy}, \textbf{Maxima}... 
    \end{itemize}
    
\item 
    Existen muchos otros entornos de cálculo numérico libres:
    \begin{itemize}
      \item {} \textbf{Python}+\textbf{NumPy}/\textbf{Scipy}, \textbf{Julia}, Scilab, Euler, Freemat...    
    \end{itemize}
    
\end{itemize}

\end{frame}


% ___________________________________________________________________________

\hypertarget{potencialidad-de-octave}{}
\pdfbookmark[0]{Potencialidad de Octave}{potencialidad-de-octave}
\begin{frame}\frametitle{Potencialidad de Octave}
\begin{itemize}
\item {} 
Curva de acceso poco empinada

\item {} 
Incluye ``de serie'':
\begin{itemize}
\item {} 
Operaciones matriciales

\item {} 
Resolución de (sistemas de) ecuaciones lineales y no lineales

\item {} 
Cálculo de autovalores y autovectores

\item {} 
Cálculo de raíces de ecuaciones no lineales

\item {} 
Aproximación numérica de funciones y datos (interpolación,
FFT, funciones spline,...)

\item {} 
Cálculo numérico de derivadas e integrales

\item {} 
Resolución numérica de ecuaciones diferenciales

\item {} 
Representación de gráficas 2D y 3D 

\item {} 
...

\end{itemize}

\item Compatibilidad con Matlab...

\end{itemize}

\end{frame}

% ___________________________________________________________________________


\hypertarget{octave-y-matlab}{}
\pdfbookmark[0]{Octave y Matlab}{octave-y-el-lenguaje-matlab}
\begin{frame}\frametitle{Octave y Matlab}
  \begin{itemize}
  \item 
    Alta compatibilidad con el ``\textbf{lenguaje Matlab}''
    \begin{itemize}
    \item {} ventajas
      \begin{itemize}
      \item fácil migración desde Matlab
      \item hereda puntos fuertes
      \end{itemize}
    \item {} inconvenientes
      \begin{itemize}
      \item freno a la innovación propia
      \item hereda puntos débiles
      \end{itemize}
    \end{itemize}
  
  \item ¿Dónde se pierde la compatibilidad?
    \begin{itemize}
    \item Extensiones al lenguaje Matlab
    \item ``Toolkits''
    \end{itemize}

  \item Filosofía propia: ``\emph{lenguaje de scripting científico}''
    \begin{itemize}
    \item Se complementa con otras herramientas libres
    \item Entorno gráfico similar a Matlab + integración con otros editores/entornos... 
    \end{itemize}

  \end{itemize}
\end{frame}

% % ___________________________________________________________________________


\hypertarget{qu-es-octave-ii}{}
\pdfbookmark[0]{Octave y Matlab}{qu-es-octave-ii}
\begin{frame}\frametitle{Lenguajes interpretados en cálculo científico}

\begin{itemize}
\item {} 
Lenguajes \emph{interpretados}:
\begin{itemize}
\item {} 
Reducción del ciclo de desarrollo (relación esfuerzo/resultado)

\item {} 
Polivalencia, acceso a herramientas de alto nivel

\item {}
Concentrarse en las dificultades del problema (no en
  las del lenguaje)

\end{itemize}

\item {} 
Lenguajes \emph{compilados}: Potencia de cálculo

\item {} 
Lenguajes interpretado como lenguajes ``\emph{pegamento}'': lo mejor de dos mundos
\begin{itemize}
\item 90\% del código: lenguaje interpertado
\item 10\% restante (zonas críticas en potencia de calculo): lenguaje
  compilado
\end{itemize}
\item Para tareas científicas:
\begin{itemize}
\item \textbf{Octave}
\item Python
\item ...
\end{itemize}
\end{itemize}

\end{frame}

% % ___________________________________________________________________________

% \hypertarget{interfaces}{}
% \pdfbookmark[0]{Interfaces}{interfaces}
% \begin{frame}\frametitle{Interfaces}
% \begin{itemize}
% \item {} 
% Octave incluye una interfaz muy sencilla en forma de
% consola (no gráfica)

% \item {} 
% Existen algunas alternativas:
% \begin{itemize}
% \item {} 
% Modo Octave \textbf{Emacs}

% \item {} 
% \textbf{qtOctave}

% \item {} 
% Otras

% \end{itemize}

% \end{itemize}

% \end{frame}


\subsection{Primeros pasos}

\hypertarget{primera-sesi-n-con-qtoctave}{}
\pdfbookmark[0]{Primera sesión con Octave}{primera-sesi-n-con-qtoctave}
\begin{frame}
\frametitle{Primera sesión con Octave}
Arrancamos Octave en modo gráfico...
\begin{itemize}

\item {} 
Distintas secciones

\begin{itemize}

\item {} 
  Explorador de archivos

\item {} Espacio de trabajo: edición de variables
\item {} Historial de comandos
\item {} 
 {Ventana de ``comandos''}: usando Octave como una calculadora

\end{itemize}

\item Otras utilidades:
\begin{itemize}
\item Editor integrado
  \item {Acceso a la ayuda de Octave}%
  \footnote{\tiny\url{chttps://docs.octave.org/interpreter/index.html}}
\item Entorno mejorado para la representación de gráficos
\item ...
\end{itemize}

\end{itemize}

\end{frame}



% ___________________________________________________________________________

\begin{frame}[containsverbatim]
\frametitle{Octave como una calculadora (I)}

\begin{itemize}
\item Operadores numéricos habituales: 
  \par \verb|+|, \verb|-|, \verb|*|, \verb|/|, \verb|^|
\item Paréntesis para modificar la preferencia
\item El resultado se almacena en la variable \verb|ans|
\item Re-editar órdenes: $\uparrow$, $\downarrow$, \emph{(Mayúsculas
    +)} $\rightarrow$, $\leftarrow$ 
\end{itemize}

\small
\begin{Octave}
>> 1+3*1.5
ans =  5.5000
>> (2+3)^10
ans =  9765625
>> ans/5^7
ans =  125
\end{Octave}


\end{frame}


\begin{frame}[containsverbatim]
\frametitle{Octave como una calculadora (II)}
Estilo Matlab...
\begin{itemize}
\item {} 
Comillas simples para introducir una cadena de texto

\item {} 
Porcentaje para comentarios

\item {} 
Tres puntos para continuación de línea

\item {} 
Punto y coma para un retorno de carro sin salida

\end{itemize}

\small
\begin{Octave}
>> 'Esto es una cadena de texto'
ans = Esto es una cadena de texto
>> 20000000 + 3100020500 +...
> 13
ans =  3.1200e+09
>> 2+2; % No se muestra el resultado
\end{Octave}

\end{frame}

\begin{frame}[containsverbatim]
\frametitle{Octave como una calculadora (II)}
Estilo Octave (no compatible!)...
\begin{itemize}
\item {} 
Comillas dobles para introducir una cadena de texto

\item {} 
Almohadilla (\verb|#|) para comentarios

\item {} 
Barra invertida (\verb|\|) para continuación de línea
\end{itemize}

Similar a otros lenguajes de scripting: python, bash...

\small
\begin{Octave}
>> ``Esto es una cadena de texto''
ans = Esto es una cadena de texto
>> 20000000 + 3100020500 +\
> 13
ans =  3.1200e+09
>> 2+2; # No se muestra el resultado
\end{Octave}

\end{frame}


\begin{frame}[containsverbatim]
\frametitle{Más madera...}
\small
\begin{Octave}
>> 1/7          % Aritmetica punto fijo
ans =  0.14286
>> format long  % Mostrar doble precisi'on
>> 1/7 
ans =  0.142857142857143
>> format       % O bien 'format short'
>> sin(pi)      %  sen(pi) = 0 ?
ans =  1.2246e-16
>> eps          % El "cero" de la maquina
ans =  2.2204e-16
>> x=1+2*i; y=2+3*i;  % Asignacion, complejos
>>  x*y
ans = -4 + 7i
\end{Octave}
\end{frame}

\begin{frame}[containsverbatim]
\frametitle{Variables}
Órdenes:
\begin{itemize}
\item \verb|who|: listado de variables
\item \verb|whos|: listado detallado
\item \verb|clear x y ...|: borrar variables \verb|x|, \verb|y|,...
\item \verb|clear all|: borrarlas todas
\item \verb|help orden|: información sobre una orden de octave
\item \verb|save fichero x y ...|\footnote{Octave puede grabar en numerosos
    formatos. Véase la ayuda}
\item \verb|load fichero x y ... | ó \verb|load fichero|
\end{itemize}
Variables reservadas:
\begin{itemize}
\item \verb|ans|, \verb|eps|, \verb|inf|, \verb|nan|, \verb|pi|,
  \verb|i|, \verb|j|, \verb|nargin|, \verb|nargout|, \verb|realmin|,
  \verb|realmax|, \verb|flintmax|, \verb|varargin|, \verb|varargout|
\end{itemize}
\end{frame}

\begin{frame}[containsverbatim]
\frametitle{Ejercicios}
\begin{enumerate}
\item
  Definir las variables:
  \begin{itemize}
  \item $\displaystyle x=e^{\pi-1}$
  \item $\displaystyle y=\frac{1}{1+\frac{1}{x+2}}$
  \item $\displaystyle z=\sqrt{x+y}$
  \item $\displaystyle t=\log(x+y+z)$
  \end{itemize}
\item Almacenarlas en un fichero (binario) llamado ``variables.dat''
\item Borrar todas las variables y volverlas a cargar desde el fichero
\end{enumerate}
\bigskip
``Para subir nota'': Grabar estas variables en formato texto (\texttt{save -ascii}) llamado ``variables.txt'':
\end{frame}

% ___________________________________________________________________________

%%% Local Variables: 
%%% mode: latex
%%% TeX-master: "modulo_octave"
%%% End: 
